\NormalFont\ShortTitle{1 Timotei}

{\MT Prima scrisoare a lui Pavel către Timotei


\par }\ChapOne{1}{\PP \VerseOne{1}Pavel, apostol al lui Isus Hristos, după porunca lui Dumnezeu, Mântuitorul nostru, și a Domnului Isus Hristos, nădejdea noastră,

\VS{2}către Timotei, adevăratul meu copil în credință: Har, îndurare și pace de la Dumnezeu Tatăl nostru și de la Hristos Isus, Domnul nostru.


\par }{\PP \VS{3}După cum v-am îndemnat când mergeam în Macedonia, rămâneți la Efes, ca să porunciți unor oameni să nu învețe o altă învățătură,

\VS{4}și să nu dați atenție la mituri și la genealogii nesfârșite, care provoacă certuri, în loc să dea atenție la administrarea lui Dumnezeu, care este în credință.

\VS{5}Dar scopul acestei porunci este dragostea dintr-o inimă curată, o conștiință bună și o credință sinceră,

\VS{6}de la care unii, ratând ținta, s-au abătut spre vorbărie deșartă,

\VS{7}dorind să fie învățători ai legii, deși nu înțeleg nici ceea ce spun, nici despre ceea ce afirmă cu tărie.


\par }{\PP \VS{8}Dar noi știm că Legea este bună, dacă o folosește cineva în mod legal,

\VS{9}știind că Legea nu a fost făcută pentru cel neprihănit, ci pentru cei fără de lege și nesupuși, pentru cei nelegiuiți și păcătoși, pentru cei necredincioși și profanatori, pentru ucigașii de tați și ucigașii de mame, pentru ucigașii de oameni,

\VS{10}pentru cei imorali din punct de vedere sexual, pentru homosexuali, pentru traficanții de sclavi, pentru mincinoși, pentru sperjurori și pentru orice alt lucru contrar bunei învățături,

\VS{11}potrivit cu Buna Vestire a slavei binecuvântatului Dumnezeu, care mi-a fost încredințată.


\par }{\PP \VS{12}Mulțumesc Celui ce m-a făcut în stare, Hristos Isus, Domnul nostru, pentru că m-a socotit credincios, punându-mă în slujbă,

\VS{13}deși eram hulitor, prigonitor și obraznic. Cu toate acestea, am obținut îndurare, pentru că am făcut-o din necredință, în necunoștință de cauză.

\VS{14}Harul Domnului nostru a abundat nespus de mult cu credința și cu dragostea care este în Cristos Isus.

\VS{15}Este credincioasă și vrednică de toată acceptarea afirmația că Isus Hristos a venit în lume ca să mântuiască pe păcătoși, dintre care eu sunt cel dintâi.

\VS{16}Totuși, din această cauză am obținut îndurare, pentru ca, în mine mai întâi, Isus Hristos să-și arate toată răbdarea ca exemplu pentru cei care urmau să creadă în el pentru viața veșnică.

\VS{17}Acum, Regelui veșnic, nemuritor și invizibil, lui Dumnezeu, care este singurul înțelept, să fie cinstea și gloria în vecii vecilor. Amin.


\par }{\PP \VS{18}Îți încredințez această învățătură, copilul meu Timotei, conform profețiilor care ți-au fost date mai înainte, pentru ca prin ele să duci lupta cea bună,

\VS{19}păstrând credința și o conștiință bună, pe care unii, după ce au îndepărtat-o, au naufragiat în ceea ce privește credința,

\VS{20}printre care se numără Himeneu și Alexandru, pe care i-am dat pe mâna Satanei, ca să fie învățați să nu hulească.



\par }\Chap{2}{\PP \VerseOne{1}Îndemn, așadar, mai întâi de toate, să se facă cereri, rugăciuni, mijlociri și mulțumiri pentru toți oamenii,

\VS{2}pentru împărați și pentru toți cei care sunt în poziții înalte, ca să ducem o viață liniștită și pașnică în toată evlavia și respectul.

\VS{3}Căci acest lucru este bun și plăcut înaintea lui Dumnezeu, Mântuitorul nostru,

\VS{4}care dorește ca toți oamenii să fie mântuiți și să ajungă la deplina cunoaștere a adevărului.

\VS{5}Căci există un singur Dumnezeu și un singur mijlocitor între Dumnezeu și oameni, omul Isus Cristos,

\VS{6}care s-a dat pe sine însuși ca răscumpărare pentru toți, mărturie la timpul potrivit,

\VS{7}pentru care am fost numit predicator și apostol — spun adevărul în Cristos, nu mint —, învățător al neamurilor în credință și adevăr.


\par }{\PP \VS{8}Așadar, doresc ca oamenii să se roage în fiecare loc, ridicând mâini sfinte, fără mânie și fără îndoială.

\VS{9}La fel, ca și femeile să se împodobească în haine decente, cu modestie și cuviinciozitate, nu cu păr împletit, cu aur, cu perle sau cu haine scumpe,

\VS{10}ci cu fapte bune, ceea ce se cuvine femeilor care mărturisesc evlavia.

\VS{11}Femeia să învețe în liniște, cu deplină supunere.

\VS{12}Dar eu nu permit ca o femeie să învețe și nici să exercite autoritate asupra unui bărbat, ci să fie în liniște.

\VS{13}Căci mai întâi a fost format Adam, apoi Eva.

\VS{14}Adam nu a fost înșelat, dar femeia, fiind înșelată, a căzut în neascultare;

\VS{15}dar ea va fi salvată prin nașterea copiilor, dacă vor continua în credință, în iubire și în sfințenie cu sobrietate.



\par }\Chap{3}{\PP \VerseOne{1}Iată o vorbă credincioasă: cine vrea să fie supraveghetor, dorește o lucrare bună.

\VS{2}Așadar, supraveghetorul trebuie să fie fără prihană, soț al unei singure soții, cumpătat, cumpătat, cumpătat, modest, ospitalier, bun la învățătură;

\VS{3}să nu fie băutor, să nu fie violent, să nu fie lacom de bani, ci blând, să nu fie certăreț, să nu fie lacom;

\VS{4}să fie unul care își conduce bine propria casă, având copiii în supunere cu toată reverența;

\VS{5}(căci cum ar putea cineva care nu știe să-și conducă propria casă să aibă grijă de adunarea lui Dumnezeu?)

\VS{6}nu un nou convertit, ca nu cumva, îngâmfat, să nu cadă în aceeași osândă ca și diavolul.

\VS{7}Mai mult, trebuie să aibă o mărturie bună din partea celor de afară, ca să nu cadă în ocară și în capcana diavolului.


\par }{\PP \VS{8}La fel și slujitorii trebuie să fie cuviincioși, să nu aibă o limbă cu două tăișuri, să nu fie dependenți de mult vin, să nu fie lacomi de bani,

\VS{9}și să aibă o conștiință curată în ceea ce privește taina credinței.

\VS{10}Să fie și ei mai întâi puși la încercare; apoi să slujească, dacă sunt ireproșabili.

\VS{11}Soțiile lor, la fel, trebuie să fie respectuoase, să nu fie defăimătoare, cumpătate și credincioase în toate lucrurile.

\VS{12}Slujitorii să fie soți ai unei singure soții, să își conducă bine copiii și casele lor.

\VS{13}Căci cei care au slujit bine își dobândesc o poziție bună și o mare îndrăzneală în credința care este în Hristos Isus.


\par }{\PP \VS{14}Vă scriu aceste lucruri, nădăjduind să vin curând la voi,

\VS{15}dar dacă voi aștepta mult, ca să știți cum trebuie să vă purtați în casa lui Dumnezeu, care este Adunarea Dumnezeului celui viu, stâlpul și temelia adevărului.

\VS{16}Fără îndoială, taina evlaviei este mare:

\par }{\Q Dumnezeu a fost revelat în carne și oase,

\par }{\Q îndreptățit în spirit,

\par }{\Q văzut de îngeri,

\par }{\Q propovăduit printre națiuni,

\par }{\Q crezut în lume,

\par }{\Q și primit în slavă.



\par }\Chap{4}{\PP \VerseOne{1}Dar Duhul spune în mod expres că, în vremurile de pe urmă, unii se vor depărta de la credință, ascultând de duhuri seducătoare și de învățături ale demonilor,

\VS{2}prin ipocrizia unor oameni care spun minciuni, care-și lovesc conștiința ca un fier încins,

\VS{3}care interzic căsătoria și poruncesc să se abțină de la mâncărurile pe care Dumnezeu le-a creat pentru a fi primite cu mulțumire de cei ce cred și cunosc adevărul.

\VS{4}Căci orice făptură a lui Dumnezeu este bună și nimic nu este de lepădat, dacă este primit cu mulțumire.

\VS{5}Căci este sfințită prin cuvântul lui Dumnezeu și prin rugăciune.


\par }{\PP \VS{6}Dacă vei învăța pe frați aceste lucruri, vei fi un bun slujitor al lui Hristos Isus, hrănit cu cuvintele credinței și cu buna învățătură pe care ai urmat-o.

\VS{7}Dar refuzați fabulele profane și cele ale babelor. Exersează-te spre evlavie.

\VS{8}Căci exercițiul trupesc are o anumită valoare, dar evlavia are valoare în toate lucrurile, având făgăduința vieții de acum și a celei viitoare.

\VS{9}Acest cuvânt este credincios și vrednic de toată acceptarea.

\VS{10}Căci în acest scop ne străduim și suferim ocara, pentru că ne-am pus încrederea în Dumnezeul cel viu, care este Mântuitorul tuturor oamenilor, mai ales al celor care cred.

\VS{11}Poruncește și învață aceste lucruri.


\par }{\PP \VS{12}Nimeni să nu disprețuiască tinerețea voastră, ci să fiți o pildă pentru cei credincioși, prin cuvânt, prin felul vostru de viață, prin dragoste, prin duh, prin credință și prin curăție.

\VS{13}Până voi veni eu, fiți atenți la citire, la îndemn și la învățătură.

\VS{14}Nu neglijați darul care este în voi, care v-a fost dat prin profeție, prin punerea mâinilor bătrânilor.

\VS{15}Fiți sârguincioși în aceste lucruri. Dăruiește-te cu totul lor, pentru ca progresul tău să fie descoperit tuturor.

\VS{16}Fiți atenți la voi înșivă și la învățătura voastră. Continuați în aceste lucruri, căci, făcând aceasta, vă veți salva atât pe voi înșivă, cât și pe cei care vă ascultă.



\par }\Chap{5}{\PP \VerseOne{1}Nu mustrați pe cel mai în vârstă, ci îndemnați-l ca pe un tată; pe cei mai tineri, ca pe niște frați;

\VS{2}pe cele mai în vârstă, ca pe niște mame; pe cele mai tinere, ca pe niște surori, în toată curăția.


\par }{\PP \VS{3}Cinstiți văduvele care sunt văduve cu adevărat.

\VS{4}Dar, dacă vreo văduvă are copii sau nepoți, să învețe mai întâi să dea dovadă de evlavie față de familia lor și să-și răsplătească părinții, căci aceasta este plăcut înaintea lui Dumnezeu.

\VS{5}Or, cea care este văduvă cu adevărat și pustie, își pune nădejdea în Dumnezeu și stăruie în cereri și rugăciuni noapte și zi.

\VS{6}Dar cea care se dăruiește plăcerilor este moartă cât trăiește.

\VS{7}De asemenea, poruncește aceste lucruri, ca să fie fără prihană.

\VS{8}Dar dacă cineva nu are grijă de ai săi și mai ales de casa sa, a tăgăduit credința și este mai rău decât un necredincios.


\par }{\PP \VS{9}Nimeni să nu fie înscrisă ca văduvă în vârstă de mai puțin de șaizeci de ani, după ce a fost soția unui singur bărbat,

\VS{10}și a fost aprobată prin fapte bune, dacă a crescut copii, dacă a fost ospitalieră cu străinii, dacă a spălat picioarele sfinților, dacă a ușurat pe cei necăjiți și dacă a urmat cu sârguință orice lucrare bună.


\par }{\PP \VS{11}Dar refuză pe văduvele mai tinere, pentru că, după ce au devenit desfrânate împotriva lui Hristos, vor să se căsătorească,

\VS{12}fiind osândite, pentru că au respins primul lor angajament.

\VS{13}În plus, ele învață și să fie leneșe, umblând din casă în casă. Nu numai leneși, ci și bârfitori și băgători de seamă, spunând lucruri pe care nu trebuie să le spună.

\VS{14}Așadar, doresc ca văduvele mai tinere să se căsătorească, să aibă copii, să conducă gospodăria și să nu dea prilej adversarului să insulte.

\VS{15}Căci deja unii s-au abătut după Satana.

\VS{16}Dacă vreun bărbat sau vreo femeie credincioasă are văduve, să le ușureze, și să nu lase adunarea să fie împovărată, ca să ușureze pe cele care sunt într-adevăr văduve.


\par }{\PP \VS{17}Să fie socotiți vrednici de o cinste dublă bătrânii care conduc bine, mai ales cei ce lucrează în cuvânt și în învățătură.

\VS{18}Căci Scriptura zice: “Să nu pui botniță boului când calcă grâul”. Și:
{\WJ{“Lucrătorul este vrednic de plata sa”.
}}


\par }{\PP \VS{19}Să nu primești o acuzație împotriva unui bătrân decât pe baza a doi sau trei martori.

\VS{20}Pe cei ce păcătuiesc, mustrați-i în fața tuturor, pentru ca și ceilalți să se teamă.

\VS{21}Vă poruncesc, în fața lui Dumnezeu, a Domnului Isus Hristos și a îngerilor aleși, să respectați aceste lucruri fără prejudecăți, fără să faceți nimic cu părtinire.

\VS{22}Să nu puneți mâinile în grabă asupra nimănui. Nu fiți părtași la păcatele altora. Păstrați-vă puritatea.


\par }{\PP \VS{23}Nu mai bea numai apă, ci bea și puțin vin, pentru sănătatea stomacului tău și pentru bolile tale frecvente.


\par }{\PP \VS{24}Păcatele unora sunt evidente, care îi preced la judecată, iar altele vin mai târziu.

\VS{25}La fel și faptele bune sunt evidente, iar cele care sunt altfel nu pot fi ascunse.



\par }\Chap{6}{\PP \VerseOne{1}Toți cei ce sunt robi sub jug să socotească pe stăpânii lor vrednici de toată cinstea, ca să nu fie hulit numele lui Dumnezeu și învățătura.

\VS{2}Cei care au stăpâni credincioși, să nu-i disprețuiască, pentru că sunt frați, ci mai degrabă să le slujească, pentru că cei care au parte de binefacere sunt credincioși și iubiți. Învățați și îndemnați aceste lucruri.


\par }{\PP \VS{3}Dacă cineva învață o altă învățătură și nu se supune cuvintelor sănătoase, cuvintelor Domnului nostru Isus Hristos și învățăturii care este potrivit cu evlavia,

\VS{4}este un îngâmfat, care nu știe nimic, ci este obsedat de certuri, de dispute și de lupte de cuvinte, de unde vin invidii, certuri, insulte, suspiciuni rele,

\VS{5}fricțiuni constante ale oamenilor cu minți corupte și lipsiți de adevăr, care cred că evlavia este un mijloc de câștig. Îndepărtează-te de astfel de oameni.


\par }{\PP \VS{6}Dar evlavia cu mulțumire este un mare câștig.

\VS{7}Căci noi nu am adus nimic în lume și cu siguranță nu putem duce nimic la îndeplinire.

\VS{8}Dar având hrană și îmbrăcăminte, ne vom mulțumi cu asta.

\VS{9}Dar cei care sunt hotărâți să se îmbogățească cad în ispită, în capcană și în multe pofte nebunești și dăunătoare, de natură să înece pe oameni în ruină și distrugere.

\VS{10}Căci iubirea de bani este o rădăcină a tot felul de rele. Unii au fost rătăciți de la credință din cauza lăcomiei lor și s-au străpuns singuri cu multe necazuri.


\par }{\PP \VS{11}Dar tu, om al lui Dumnezeu, să fugi de aceste lucruri și să urmezi dreptatea, evlavia, credința, dragostea, perseverența și blândețea.

\VS{12}Luptă lupta cea bună a credinței. Apucă-te de viața veșnică la care ai fost chemat și ai mărturisit buna mărturisire în fața multor martori.

\VS{13}Vă poruncesc, înaintea lui Dumnezeu, care dă viață tuturor lucrurilor, și înaintea lui Isus Cristos, care în fața lui Ponțiu Pilat a mărturisit buna mărturisire,

\VS{14}să respectați porunca fără pată, fără vină, până la arătarea Domnului nostru Isus Cristos,

\VS{15}pe care, la timpul potrivit, o va arăta el, care este binecuvântatul și singurul Conducător, Regele regilor și Domnul domnilor.

\VS{16}El singur are nemurirea, care locuiește într-o lumină de neatins, pe care nimeni nu l-a văzut și nici nu poate să-l vadă, căruia i se cuvine cinstea și puterea veșnică. Amin.


\par }{\PP \VS{17}Îndemnați-i pe cei ce sunt bogați în veacul acesta să nu se îngâmfe și să nu-și pună nădejdea în bogăția nesigură, ci în Dumnezeul cel viu, care ne dăruiește din belșug toate cele de care ne putem bucura;

\VS{18}să facă binele, să fie bogați în fapte bune, să fie gata să împartă, gata să împartă;

\VS{19}să-și facă rezerve pentru ei înșiși o temelie bună pentru vremurile viitoare, ca să poată dobândi viața veșnică.


\par }{\PP \VS{20}Timotei, păzește ceea ce ți s-a încredințat, ferindu-te de vorbăria deșartă și de opozițiile celor ce se numesc în mod greșit știință,

\VS{21}pe care le fac unii, și care s-au abătut de la credință.

\par }{\PP Harul să fie cu voi. Amin.


\par }